% Category codes are what a character MEANS, and a document can change them at
% any point -- this is the reconfigurability that makes TeX's mouth a machine
% rather than a fixed grammar.
%
% INITEX starts with `{` and `}` ordinary, which is why every example here opens
% by assigning them; plain.tex does the same thing, it is just already done for
% you by the time a normal document runs.
\catcode`\{=1 \catcode`\}=2 \catcode`\#=6
\def\show#1{[#1]}
\message{\show{grouped}}
% A character can be made a letter, which lets it appear in a control word --
% the trick plain.tex uses for `@` so internal macros cannot be typed by
% accident.
\catcode`\@=11
\def\intern@l{PRIVATE}
\message{\intern@l}
\catcode`\@=12
% And back to ordinary: now `@` is just a character again.
\message{at-sign is \string@}
\end
